\reviewsection{Meaningful Communication Is Authorship, Not Performance}

At the beginning of Module 3, I described assistive technology as a way to reduce barriers and AAC as a multimodal communication system rather than a single device. Gesture, gaze, objects, images, text, signs, low-tech boards, devices, and speech can all belong to that system. Communication access should not depend on age, a test score, diagnosis, or proving readiness. The module did not replace that understanding. It made the standard more demanding. My Module 1 journal, \textit{Who Gets to Name, Interpret, and Design?}, asked how labels, lived experience, and educational environments shape what the world recognizes as developmental difference \citep{garciabautista2026module1journal}. My Module 2 journal then asked who gains power after support. Module 3 joins those questions: who is recognized as the author of a message, which communication is treated as credible, and can that message alter the adult's interpretation or the environment?

Meaningful communication is not defined by speech, speed, eye contact, independent device use, or whether a response matches an adult's expectation. It exists when a student's message is recognized as theirs, understood or respectfully repaired, and able to influence learning, relationships, support, and what happens next. AAC can make that authorship more accessible, but device presence does not prove communication power. There is no age, cognitive score, or developmental milestone that a student must reach before communication support can help \citep{ashaAAC}. \citet{norrie2021context} make the difference between device presence and communication access visible. In their five-month study of one special-education school, only six of approximately 180 pupils with complex communication needs were active high-tech AAC users. Their findings locate barriers not only in a device or learner but in training, technical support, staffing, policy, infrastructure, and the fit between technology and daily life. The study changes the question from ``Why is the student not using the device?'' to ``What conditions make this communication system usable across real activities and partners?''

That distinction matters because a learner can technically possess AAC while having access to only a narrow adult agenda. A system limited to requesting food, answering teacher questions, or following directions does not provide the language needed to disagree, joke, ask an unexpected question, express identity, repair a misunderstanding, participate academically, or refuse. Meaningful communication must include functions that adults did not script in advance. \citet{downing2015classroom} move that responsibility into the general education classroom by asking educators to examine actual routines, exchanges, peer relationships, and communication opportunities. A student can be physically present and still have little access to authorship. Predictable routines may help a learner anticipate an exchange, but meaningful participation also requires vocabulary and partner support for unpredictable comments, questions, disagreement, humor, and repair. Communication cannot be reduced to selecting between two adult-approved choices.
